% =====================================================================
%  Section 1 -- Introduction
% =====================================================================

\section{Introduction}
\label{sec:intro}

Floods are the most frequent and economically damaging natural disasters globally, accounting for 38.7\% of natural disasters and affecting 43\% of disaster-affected populations between 2000 and 2009~\cite{du2010health}. Rapid urbanisation continues to amplify exposure~\cite{hemmati2020urban}. Because traditional hydrological and hydraulic models suffer from high computational demand, coarse spatial resolution, and limited real-time adaptability~\cite{kumar2023flood}, deep-learning segmentation has emerged as a leading approach to automated flood-extent delineation from satellite, drone, and ground imagery~\cite{bentivoglio2022deep}.

Recent benchmark studies in this space have converged on encoder--decoder architectures --- UNet~\cite{ronneberger2015unet}, SegNet~\cite{badrinarayanan2017segnet}, and DeepLabV3$+$~\cite{chen2018encoder} --- paired with strong feature-extraction backbones, most commonly EfficientNet~\cite{tan2019efficientnet} and ResNet~\cite{he2016deep}. Karc{\i} \emph{et al.}~\cite{karci2026deep} provide the most comprehensive systematic comparison of these combinations on the public Flood Semantic Segmentation Dataset (FSSD), reporting that UNet $+$ EfficientNet achieves the highest segmentation performance (IoU 0.927, F1 0.957, accuracy 0.971) and conclude that this combination ``is ideal for near real-time operational deployment.''

A careful reading of the broader flood-segmentation literature, however, reveals that this operational-deployment claim has not actually been evaluated. Across the 26 related works surveyed by Karc{\i} \emph{et al.}~\cite{karci2026deep} and adjacent studies on SAR~\cite{pechmay2023flooded} and multi-modal flood mapping~\cite{chamatidis2024vision}, nearly every paper reports segmentation accuracy alone. The exception, Hern\'andez \emph{et al.}~\cite{hernandez2022flood}, times PSPNet/DeepLabV3$+$/U-Net on CUDA-equipped NVIDIA Jetson boards --- where the largest configurations fail to run on the Jetson Nano at all --- but \emph{no study reports latency on commodity CPU- or browser-class hardware, and none reports post-quantization accuracy or quantization stability.} This omission is consequential. UNet $+$ EfficientNet-B7 --- the configuration favoured in~\cite{karci2026deep} --- has approximately 63\,M parameters and 50\,GFLOPs at $256{\times}256$ input resolution. Such a model cannot run on the drones, embedded sensors, Raspberry-Pi-class community stations, or offline tablet apps that humanitarian organisations such as UN-SPIDER and the ICRC actually deploy in flood zones, where reliable internet connectivity is the exception rather than the rule~\cite{munawar2021review}.

\subsection{Contributions}
\label{sec:contributions}

This paper makes four contributions.

\begin{enumerate}
\item \textbf{First systematic CPU-and-browser edge-deployment benchmark for flood segmentation.} We report FP32 and INT8 inference latency for three lightweight UNet students on two deployment targets --- Apple M2~Pro via ONNX Runtime CPU, and browser WebAssembly SIMD via ONNX Runtime Web 1.17 --- together with per-model parameter counts, FLOPs at $256{\times}256$, ONNX disk size, and FP32-to-INT8 size-reduction factors. ONNX serves as the deployment hub: every measured latency reflects the same exported artifact running under a different execution provider. We benchmark these two commodity software runtimes on a single Apple~M2~Pro host; measurement on physically constrained edge devices (Raspberry~Pi, mobile SoCs, embedded ARM) is out of scope for this study and deferred to future work (Section~\ref{sec:limitations}). Prior and concurrent edge-oriented flood-segmentation work targets CUDA-equipped Jetson boards~\cite{hernandez2022flood,schmalstieg2026distillation}; we characterise the GPU-free deployment envelope that field laptops, tablets, and web browsers actually provide.

\item \textbf{When knowledge distillation helps flood segmentation, and what it transmits.} On FSSD, across three architectures $\times$ four KD configurations $\times$ three folds, no KD configuration improves on the no-KD baseline: ImageNet-pretrained lightweight students match or exceed the UNet$+$EfficientNet-B0 teacher without distillation. A per-image paired bootstrap over the pooled held-out set ($n=399$ images) confirms this and shows response-based KD is significantly \emph{harmful} on two of the three students. We attribute the result to the saturated 663-image benchmark leaving the teacher no headroom. On a less saturated combined corpus, where the teacher does out-segment the students, combined KD raises flood-only IoU significantly for all three students. It also passes on the teacher's false alarms on water-free scenes. Distillation transfers the teacher's errors along with its skill, so the teacher has to be audited for the deployment-relevant failure first.

\item \textbf{Architecture-dependent post-training quantization stability characterisation.} Under an identical ONNX static QDQ INT8 recipe (Percentile 99.999\% calibration, per-tensor scales for both weights and activations, quantization restricted to Conv/Gemm/MatMul operators), EfficientNet-Lite0 quantizes cleanly ($\Delta \mathrm{IoU} = -0.043 \pm 0.016$) while MobileNetV3-Small and MobileViT-XXS suffer 0.45--0.46 IoU collapses with cross-fold standard deviations of 0.19--0.34. We trace the failure modes to (i)~the HardSwish activations in MobileNetV3-Small, whose measured post-activation range is $12$--$17\times$ wider than EfficientNet-Lite0's bounded ReLU6 (Section~\ref{sec:results-quant}), and (ii)~the LayerNorm and self-attention blocks in MobileViT-XXS, whose severe inter-channel variation and non-uniform attention distributions exceed what per-tensor static calibration can represent~\cite{lin2022fqvit}. The implication for practitioners: \emph{bounded-activation backbones (ReLU6) post-training-quantize reliably for flood segmentation; HardSwish and attention-based backbones require quantization-aware training (QAT).}

\item \textbf{The false-alarm failure of flood-only training data, and its remedy.} Beyond the cross-sensor Sen1Floods11 collapse (Section~\ref{sec:results-ood}), we evaluate every FSSD-trained model, FP32 and INT8, on two \emph{within-regime} RGB test sets whose disjointness from FSSD we verify by perceptual-hash deduplication: AIFloodSense (94 ground-level images) and the held-out FloodNet test split (424 UAV images, 215 water-free). FP32 students transfer as well as the teacher (0.69--0.72 IoU), but every model false-alarms on 72--100\% of water-free scenes, a consequence of flood-only training data that no FSSD-internal metric reveals. INT8 degradation also grows under domain shift, up to outright collapse of the architectures that already quantized marginally in-domain. Retraining on a combined corpus that includes water-free images brings student false-alarm rates down to 3--9\% and flood-only IoU on FloodNet up to 0.70--0.78, which turns the failure into a concrete data requirement for unattended flood monitoring.
\end{enumerate}

We make our trained checkpoints, the ONNX FP32 and INT8 exports, the per-platform benchmarking harnesses, and a Zenodo-archived release of the full code repository publicly available.

\subsection{Scope and pivot from initial framing}
\label{sec:scope-pivot}

This work was originally framed as a knowledge-distillation recipe for compressing a heavy flood-segmentation teacher. The empirical results did not support that framing, and we report the negative-KD finding rather than search for a configuration that would justify the original framing. The contribution pivots to the per-platform edge benchmark, the architecture-dependent PTQ characterisation, and the conditions under which distillation pays off. Testing the negative result on a second, less saturated corpus showed that it is a property of the benchmark rather than of distillation, and that where distillation does help it carries a cost the headline IoU does not show. We hope this is more useful to practitioners than another KD recipe.

The remainder of the paper is organised as follows. Section~\ref{sec:related} reviews relevant prior work on flood segmentation, lightweight segmentation backbones, knowledge distillation, and post-training quantization. Section~\ref{sec:methods} details the FloodLite methodology. Section~\ref{sec:experimental-setup} describes the experimental setup, including the 3-fold protocol and per-platform latency rigging. Section~\ref{sec:results} reports results across Tables~\ref{tab:t1-footprints}--\ref{tab:t10-combined} and Figures~\ref{fig:f2-pareto}--\ref{fig:f4-throughput}. Section~\ref{sec:discussion} discusses the negative-KD finding, per-platform deployment recommendations, and the architecture-dependent PTQ insight. Section~\ref{sec:limitations} lists limitations. Section~\ref{sec:conclusion} concludes.
